% Tipo de documento. En este caso es un art�culo, para folios A4, tama�o de la fuente 11pt y con p�gina separada para el t�tulo
\documentclass[a4paper,11pt,titlepage]{article}

% Carga de paquetes necesarios. OrdenesArticle es un paquete personalizado
\usepackage[spanish]{babel} 
%quitado temporalmente -> \RequirePackage[T1]{fontenc}
\RequirePackage[ansinew]{inputenx} 
\usepackage[spanish,cap,cont,title,fancy]{OrdenesArticle}
\usepackage{array}
\usepackage{graphicx}
\usepackage{hyperref}
\usepackage{pifont}
\usepackage{listings}
\usepackage[usenames,dvipsnames]{color}
\usepackage{colortbl}
\usepackage{makeidx}
\hypersetup{bookmarksopen,bookmarksopenlevel=3,linktocpage,colorlinks,urlcolor=blue,citecolor=blue,
						linkcolor=blue,filecolor=blue,pdfnewwindow,
						pdftitle={Trabajo Final}, % Cambiar el t�tulo por lo que se quiera poner. Lo mismo para el autor y la materia
						pdfauthor={Juan Andrada Romero, Juan Gallardo Casero, Jose Domingo L�pez L�pez},
						pdfsubject={Pr�ctica Primer Cuatrimestre Ingenier�a Software II}}


% Macro para definir una lista personalizada 
\newenvironment{milista}%
{\begin{list}{\textbullet}%
{\settowidth{\labelwidth}{\textbullet} \setlength{\leftmargin}{\dimexpr\labelsep+\labelwidth+5pt}
\setlength{\itemsep}{\dimexpr 0.5ex plus 0.25ex minus 0.25ex}
\setlength{\parsep}{\itemsep}
\setlength{\partopsep}{\itemsep}
\addtolength{\topsep}{-7.5pt}
}}%
{\end{list}}


% Definicion del "listings" para el lenguaje Java
\lstdefinelanguage{Java}
{
 morecomment = [l]{//}, 
 morecomment = [l]{///},
 morecomment = [s]{/*}{*/},
 morestring = [b]", 
 sensitive = true,
 morekeywords = {package, static, while, switch, break, line, void, String, Object, int, Integer, instanceof, else, if, for, private, return, new, public, class, import, int, boolean, true, false, extends, final, super, protected, abstract, this, do, float, double, null}
}

% Configuraci�n del c�digo incluido para el lenguaje Java
\lstset{
  language=Java,
  basicstyle=\footnotesize,
  backgroundcolor=\color{white},
  showspaces=false,
  showstringspaces=false,
  showtabs=false,
  frame=single,
  tabsize=2,
  captionpos=b,
  breaklines=true,
  breakatwhitespace=false,
  escapeinside={\%},
  keywordstyle = \color [rgb]{0,0,1},
  commentstyle = \color [rgb]{0.133,0.545,0.133},
  stringstyle = \color [rgb]{0.627,0.126,0.941}
}


\begin{document}

% En las p�ginas de portada e �ndices, no hay encabezado ni pie de p�gina
\pagestyle{empty} 

% Se incluye la portada
\thispagestyle{empty}
\begin{center}
  {\LARGE UNIVERSIDAD DE CASTILLA-LA MANCHA} \\
  \bigskip
  {\Large ESCUELA SUPERIOR DE INFORM�TICA} \\
  \vspace{28mm}
  \includegraphics[scale=0.45, keepaspectratio]{./imagenes/esi_bw.png} \\
  \vspace{30mm}
  {\LARGE \textbf{ Ingenier�a del Software II}} \\
  \vspace{10mm}
  {\large \textsf{\textsc{SSCA - Sistema de Salud de una Comunidad}}} \\
  {\large \textsf{\textsc{Aut�noma}}} \\
  \vspace{20mm}
  {\large Juan Andrada Romero} \\
  {\large Juan Gallardo Casero} \\
  {\large Jose Domingo L�pez L�pez} \\
  \vspace{9mm}
  {\large \today}
\end{center}
\clearpage

% Texto del reverso de la portada
%%\mbox{}
%%\vspace{18cm}
%%\begin{small}
% Se ajusta la separaci�n entre p�rrafos
%%\parskip=10pt 

%%\copyright~ 2008/2009 Juan Andrada Romero. Universidad de Castilla La Mancha, Escuela Superior de Inform�tica de Ciudad Real.

%%Se permite la modificaci�n, copia y distribuci�n de este documento, seg�n la licencia de documentaci�n GNU (\url{http://www.gnu.org}).

%%Este documento fue compuesto con \LaTeX{}. Im�genes generadas con Power Point y Gimp.
%%\end{small}

%%\newpage

% En las p�ginas de �ndices y prefacio, se utiliza numeraci�n romana
\pagenumbering{Roman}

% Se crea el �ndice
\tableofcontents
% Se pasa p�gina y se a�ade el �ndice de figuras al �ndice principal
%\clearpage\phantomsection
%\addcontentsline{toc}{section}{\listfigurename}
% Se crea el �ndice de figuras
%\listoffigures
% Si se quiere crear un �ndice de tablas se pondr�a: \listoftables

\newpage

% Se ajusta la separaci�n entre p�rrafos
\parskip=10pt

% Se a�ade el prefacio al �ndice
\clearpage\phantomsection
\addcontentsline{toc}{section}{Prefacio}

% Comienza el contenido del documento. Se utilizan n�meros ar�bigos y el encabezado y pie de p�gina personalizado
\pagenumbering{arabic}
\pagestyle{fancy}

% Aqui se incluyen los archivos .tex que forman el documento
\section{Fase de Inicio}

\subsection{Especificaci�n inicial de requisitos} \label{requisitosIniciales}

Se desea desarrollar un sistema para la gesti�n de citas del sistema de salud
de una comunidad aut�noma. En la primera fase del proyecto (correspondiente al
primer cuatrimestre del curso) se debe construir la aplicaci�n que se
utilizar� en las unidades de citaci�n de los hospitales y centros de salud. En
la segunda fase (segundo cuatrimestre) se construir� una aplicaci�n web para
que pueda ser usada por parte de los beneficiarios del sistema de salud. Los
requisitos de esta segunda fase se dar�n oportunamente.

El sistema de salud de la comunidad aut�noma (SSCA) tiene un conjunto muy
amplio de beneficiarios, de cada uno de los cuales se almacenan sus datos
personales habituales (n�mero de seguridad social, NIF o NIE, nombre,
apellidos, domicilio, tel�fono fijo y m�vil, correo electr�nico, etc.).
Cada beneficiario tiene asignado un m�dico de cabecera, que pasa consulta en
un consultorio, centro de salud u hospital. De los m�dicos tambi�n se
almacenan sus datos personales.

Cuando un beneficiario desea pedir cita, se acerca a una unidad de citaci�n o
llama por tel�fono, de manera que le atiende un empleado que es quien que
maneja el sistema que se debe desarrollar en la primera fase. El beneficiario
se identifica con su n�mero de seguridad social, NIF o NIE, y entonces el
empleado ve cu�l es el m�dico de cabecera que tiene asignado. El empleado
propone al beneficiario alguno de los d�as y horas en los que el m�dico le
puede atender. Cuando el beneficiario acepta una cita, el empleado la marca y
�sta le queda reservada.

Cada m�dico puede tener un horario de consultas fijo (por ejemplo, todos los
d�as de 9 a 14) o variable (lunes y martes de 9 a 14, mi�rcoles de 12 a 14 y
de 17 a 20, ...), si bien hay d�as en que no pasan consulta por alg�n motivo
(vacaciones, quir�fano, asuntos propios, baja, etc.). Cuando un m�dico no va
pasar consulta en una fecha, el SSCA designa normalmente un sustituto. Si el
beneficiario del SSCA pide cita y el m�dico va a estar ausente, el empleado le
debe informar del nombre del m�dico que le atender� antes de darle la cita.

El sistema debe permitir la modificaci�n y anulaci�n de citas, as� como dar
citas a beneficiarios que no est�n registrados previamente. El m�dico de
cabecera de los menores de 14 a�os es un pediatra.

En ocasiones, el m�dico de cabecera (o el pediatra) puede enviar al
beneficiario a una consulta con el especialista. En este caso, el m�dico
utiliza el mismo sistema que usted tiene que desarrollar para emitir un
``volante''. El volante permite al beneficiario pedir cita para el
especialista. El m�dico de cabecera solicitar� el volante al sistema, que le
devolver� un identificador �nico de volante. Con este n�mero, el beneficiario
puede ir a la unidad de citaci�n o llamar por tel�fono y pedir cita para el
especialista. El volante relaciona, al menos, al m�dico de cabecera (o pediatra), al m�dico especialista y al beneficiario. Luego, el m�dico
especialista puede expedir m�s volantes para que el beneficiario vuelva a
acudir a su consulta o a la consulta de otro especialista.

Debido a la informaci�n confidencial que se maneja, es necesario que quede un
registro de todas las operaciones realizadas por cada usuario del sistema.

Los empleados son los usuarios del sistema y, por tanto, utilizan los equipos
inform�ticos de las unidades de citaci�n (aplicaciones cliente). Si el usuario
es administrador, adem�s de las operaciones de los usuarios con rol de
``citador'', el usuario puede:
\begin{milista}
\item Crear, modificar y eliminar otros usuarios.
\item Asignar o desasignar roles.
\item Modificar los calendarios y horarios de disponibilidad de los
m�dicos, estableciendo los m�dicos sustitutos cuando sea necesario.
\end{milista}

En modo ``citador'', el cliente mostrar� un calendario con las citas de cada
m�dico en cada d�a. Los datos de usuarios, m�dicos, citas, etc. residen tanto
en la BD principal como en la de respaldo, que son copias id�nticas una de
otra. Cuando se confirma una cita, el servidor A recibe la petici�n de alguno
de los clientes:
\begin{milista}
\item A env�a la petici�n a B
\item B la almacena en su BD de respaldo
\item B env�a a A un mensaje de confirmaci�n
\item A almacena la cita
\item A env�a un mensaje al cliente indic�ndole que la cita se ha dado
\item A actualiza el estado de los clientes: es decir, pasa a mostrar la
hora como ocupada (o como libre, si se trata de una cancelaci�n de cita)
\end{milista}

Si B falla, la transacci�n de que se trate (alta, modificaci�n o lo que sea)
no puede realizarse.

\subsection{Especificaci�n detallada de requisitos}

Para facilitar el an�lisis y el posterior dise�o e implementaci�n del sistema,
desde la fase de inicio del desarrollo su funcionalidad se dividi� en cinco
partes b�sicas, correspondientes a la gesti�n de beneficiarios, usuarios,
citas, volantes y sustituciones. En los pr�ximos apartados se explican los
requisitos detallados de cada parte del sistema.

\subsubsection{Gesti�n de beneficiarios}

El sistema permitir� la creaci�n, consulta, modificaci�n y tambi�n la
eliminaci�n de beneficiarios, a pesar de que en la especificaci�n inicial no
se hac�a referencia a ella. De cada beneficiario, el servidor almacenar� los
siguientes datos:
\begin{milista}
\item \textbf{Datos personales}: NIF, NSS (n�mero de seguridad social), nombre,
apellidos y fecha de nacimiento.
\item \textbf{Datos de contacto}: Direcci�n completa (domicilio, n�mero, piso,
puerta, localidad, provincia y c�digo postal), correo electr�nico, tel�fono
fijo y tel�fono m�vil.
\item \textbf{Datos del sistema}: Centro de salud y m�dico de cabecera o
pediatra asignado.
\end{milista}

Al registrarse en el sistema, el citador pedir� al beneficiario que elija el
centro de salud en el que quiere pasar las citas con su m�dico de cabecera o
pediatra. De esta manera, el beneficiario podr� elegir el centro de salud m�s
cercano a su ubicaci�n, o cualquiera que sea de su preferencia, para ir a ver
al m�dico de cabecera.

Una vez que el beneficiario se haya registrado, el sistema le asignar� un
m�dico de cabecera o pediatra (en funci�n de su edad) del centro seleccionado,
y el citador deber� informarle de cu�l ha sido el m�dico elegido por el
sistema.

El beneficiario podr� pedir al empleado del sistema desarrollado la consulta,
modificaci�n y eliminaci�n de sus datos registrados en el sistema, teni�ndose
que identificar para ello con su NIF o su NSS. En general, todas las
operaciones relacionadas con la gesti�n de beneficiarios las pueden realizar
los usuarios del sistema con rol de ``citador'', no siendo necesaria la
intervenci�n de los administradores. Sin embargo, los m�dicos tambi�n deber�an
poder consultar informaci�n de los beneficiarios, por si necesitaran saber su
fecha de nacimiento o alg�n dato para ponerse en contacto con ellos.

El m�dico de cabecera o pediatra asignado a cada beneficiario no se puede
modificar, salvo si el beneficiario solicita que se le cambie el centro de
salud en el que quiere pasar las citas. En este caso excepcional, que se puede
dar si el beneficiario cambia de direcci�n, el sistema le asignar� un m�dico
diferente del nuevo centro de salud seleccionado. El NIF y el NSS de un
beneficiario es fijo y no se puede modificar en ning�n caso.

Los m�dicos asignados a los beneficiarios de menos de 14 a�os son pediatras,
mientras que los asignados a beneficiarios con 14 o m�s a�os son m�dicos de
cabecera. El sistema, de forma autom�tica, tendr� que cambiar el m�dico
asociado a un beneficiario cuando cumpla los 14 a�os.

\subsubsection{Gesti�n de usuarios}

El sistema permitir� la creaci�n, consulta, modificaci�n y eliminaci�n de 
usuarios. Se considerar� que existen tres roles o tipos de usuarios:
``citadores'', administradores y m�dicos, pudiendo haber m�s de un
administrador en el sistema. De cada uno de estos usuarios, el servidor
almacenar� los siguientes datos:
\begin{milista}
\item \textbf{Datos personales}: NIF, nombre y apellidos.
\item \textbf{Datos de contacto}: Correo electr�nico, tel�fono fijo y tel�fono
m�vil.
\item \textbf{Datos del sistema}: Nombre de usuario, contrase�a y centro de
salud asignado.
\end{milista}

Adem�s, cada m�dico tiene un horario o calendario de trabajo en el que se
define a qu� horas puede pasar consulta en cada d�a de la semana (de lunes a
viernes). Por simplicidad, y debido a que no supone una limitaci�n importante,
cada horario se define con intervalos de una hora, de forma que un m�dico puede
pasar consulta de 10:00 a 14:00, o de 10:00 a 11:00 y de 12:00 a 14:00, pero
no de 10:00 a 13:30, o de 10:15 a 13:15.

La jornada laboral v�lida para los m�dicos se considera que empieza a las 9:00
de la ma�ana y finaliza a las 21:00 de la tarde. Para permitir el registro de
m�dicos de diferentes tipos y jornadas, no se establecer� una limitaci�n
inferior ni superior en el n�mero de horas de trabajo diarias ni semanales de
cada m�dico.

Los m�dicos podr�n ser de tres clases: m�dicos de cabecera, pediatras o
especialistas. De cada especialista se guardar�n, adem�s de sus datos
habituales, su especialidad m�dica, que se utilizar� para facilitar la emisi�n
de volantes a los beneficiarios.

Cuando se registre un nuevo usuario en el sistema, independientemente de su
tipo, se le asignar� de forma autom�tica el centro de salud en el que deber�
trabajar. Este centro, del que el nuevo usuario deber� ser informado, no se
podr� modificar en ning�n caso.

Por simplicidad, cada usuario s�lo podr� tener asociado un rol (y, en el caso
de los m�dicos, una clase de m�dico), que no se podr� modificar una vez que
el usuario se ha registrado, puesto que es improbable que un m�dico de cabecera
pase a ser administrador del sistema o viceversa. El �nico cambio de rol que
se deber� poder hacer es pasar el rol de un ``citador'' a administrador y
viceversa, lo cual es m�s probable que ocurra en la realidad. Por otro lado,
el NIF de un usuario es fijo y tampoco se puede cambiar.

Los usuarios podr�n pedir al administrador del sistema desarrollado la
consulta, modificaci�n y eliminaci�n de sus datos registrados en el sistema,
teni�ndose que identificar para ello con su NIF. De los m�dicos tambi�n se
podr�n cambiar, si es necesario, su horario de trabajo, para a�adir o eliminar
horas. En general, debido a su relevancia, todas las operaciones relacionadas
con la gesti�n de usuarios s�lo las pueden realizar los usuarios del sistema
con rol de administrador.

\subsubsection{Gesti�n de citas}

El sistema permitir� el registro, consulta y anulaci�n de citas a
beneficiarios. Puesto que la modificaci�n de una cita no es m�s que una
anulaci�n de la cita antigua y el registro de la nueva, y dado que el proceso
de tramitaci�n de una cita es mucho menos costoso que el de un beneficiario o
usuario, se acord� que no era necesario que el sistema permitiera la
modificaci�n directa de las citas.

Para facilitar la gesti�n y tramitaci�n de las citas, adem�s de considerar que
el horario de consultas de un m�dico se encuentra dividido en intervalos de
una hora (como ya se coment� en el anterior apartado de gesti�n de usuarios),
las citas tendr�n una duraci�n fija de 10 minutos cada una, y s�lo se podr�n
pedir en intervalos de 10 minutos. Esto significa que se podr�n pedir cita a
las 14:10 o las 19:40, pero no a las 16:22 o 20:03.

Cada m�dico podr� tener, como m�ximo, una cita asignada para un mismo d�a y
hora, mientras que, en principio, cada beneficiario podr� tener varias citas
el mismo d�a a la misma hora. No obstante, los beneficiarios s�lo podr�n pedir
cita para el m�dico de cabecera o pediatra que le asign� el sistema, o para un
especialista si disponen de un volante, por lo que, en condiciones normales,
cada beneficiario s�lo podr� pedir una �nica cita en una misma hora.

Cuando el empleado del sistema vaya a registrar una cita para un beneficiario,
se le debe mostrar a priori qui�n es el m�dico que va a hacer realmente la
consulta, teniendo en cuenta las sustituciones que le pudieran estar haciendo
al m�dico esperado por el beneficiario. Por otro lado, si el beneficiario que
quiere pedir cita no se encuentra registrado, el sistema deber� ofrecer la
posibilidad de registrar el beneficiario siguiendo el proceso que se explic�
en el apartado de gesti�n de beneficiarios.

La unidad de citaci�n desarrollada deber� permitir la consulta de las citas
registradas tanto para un cierto beneficiario como para un cierto m�dico del
sistema. Todas las operaciones relacionadas con la gesti�n de citas las podr�n
llevar a cabo de la misma forma los usuarios con rol de ``citador'' y
administrador.

\subsubsection{Gesti�n de volantes}

El sistema permitir� �nicamente la solicitud de nuevos volantes, no siendo
posible la modificaci�n, anulaci�n o consulta de volantes existentes (salvo
cuando un beneficiario va a pedir una cita utilizando un volante).

Cuando un m�dico desee emitir un volante, deber� identificar al beneficiario
que podr� utilizarlo y, a continuaci�n, el sistema le mostrar� una lista de
los m�dicos especialistas registrados en el sistema para que seleccione el
receptor del nuevo volante. Para facilitar la selecci�n del especialista, los
m�dicos se tendr�n que agrupar seg�n su especialidad m�dica indicada al
registrarse.

Antes de emitir un volante, el m�dico deber� informar al beneficiario del
centro de salud en el que pasa consulta el especialista, debido a que, al
contrario de lo que ocurre con el m�dico de cabecera o pediatra, los
especialistas a los que se env�an los beneficiarios no tienen por qu� trabajar
en el mismo centro de salud que indicaron al registrarse.

Cada volante s�lo se podr� utilizar una �nica vez para pedir cita con el
especialista, aunque si, por cualquier motivo, se anula la cita con la que un
beneficiario utiliz� un volante, podr� usarlo de nuevo para pedir cita con el
especialista en otro momento. No obstante, en cualquier caso todos los
volantes tienen una fecha de caducidad de 6 meses desde el momento de su
emisi�n, y no se podr�n utilizar a partir de esa fecha (aunque se puedan usar
para pedir una cita en un fecha posterior a la caducidad del volante).

La emisi�n de volantes es responsabilidad exclusiva de los m�dicos del sistema,
por lo que ni los citadores ni los administradores del sistema tendr�n permiso
para solicitar nuevos volantes para los beneficiarios.

\subsubsection{Gesti�n de sustituciones}

El sistema permitir� s�lo la asignaci�n de nuevas sustituciones, ya que entre
los requisitos iniciales del sistema no se encuentra la modificaci�n, anulaci�n
o consulta de las sustituciones existentes.

De la misma forma que las horas de consulta de los m�dicos se dividen en
intervalos de una hora, las sustituciones s�lo se podr�n hacer en horas
completas, pudi�ndose solicitar una sustituci�n en varias horas consecutivas o
separadas para un mismo m�dico en un mismo d�a.

Cuando un m�dico est� ausente, su horario de consultas seguir� siendo el mismo,
pero las citas las pasar� el m�dico sustituto. Por tanto, las horas que se
deber�n sustituir ser�n �nicamente aquellas en las que el m�dico sustituido
tendr�a consulta. Un m�dico no podr� sustituir a dos m�dicos diferentes en
las mismas horas, pero s� podr� sustituir a varios m�dicos el mismo d�a si no
se solapan sus horarios de consulta. Por ejemplo, un m�dico que no trabaje los
viernes podr�a sustituir en un viernes a la vez a un m�dico que trabaja de
9:00 a 14:00 y a otro que trabaja de 16:00 a 20:00.

Los administradores ser�n los �nicos usuarios que puedan asignar sustitutos a
los m�dicos. Por tanto, cuando un m�dico no vaya a poder pasar consulta en una
fecha, deber� avisar al administrador, el cual identificar� al m�dico en el
sistema y, a continuaci�n, elegir� el d�a y las horas en las que el m�dico va
a estar ausente. Si el m�dico ya estuviera siendo sustituido o estuviera
haciendo una sustituci�n en las horas elegidas, no se podr�a asignar un nuevo
sustituto.

Con el fin de hacer m�s sencilla para el administrador la tarea de asignar
sustitutos a los m�dicos, el sistema deber� generar autom�ticamente una lista
con todos los m�dicos del sistema que puedan hacer una sustituci�n pedida. Un
m�dico podr� ser elegible para realizar una sustituci�n si cumple los
siguientes requisitos:
\begin{milista}
\item Ser exactamente de la misma categor�a que el m�dico sustituido, lo cual
incluye la especialidad m�dica en el caso de que se vaya a sustituir un
especialista.
\item Trabajar en el mismo centro de salud que el m�dico sustituido, para que
el beneficiario no se tenga que desplazar por culpa de la sustituci�n.
\item No tener horario de consultas en alguna de las horas que se tienen que
sustituir.
\item No estar pendiente de realizar una sustituci�n en alguna de las horas que
se van a tener que sustituir, incluso si a�n no tiene ninguna cita asignada en
esas horas.
\end{milista}

\subsubsection{Otras consideraciones}

Para permitir que el sistema se pueda ejecutar sobre un entorno distribuido,
la unidad de citaci�n, el servidor front-end, el servidor de respaldo, la base
de datos principal y la base de datos secundaria deben poder estar en m�quinas
diferentes sin que ello afecte al funcionamiento del sistema.

Todas las operaciones que realicen los usuarios del sistema deben quedar
registradas en la base de datos por si fuera necesaria consultarlas m�s
adelante. Adem�s, el servidor front-end debe mostrar por pantalla los
resultados de las �ltimas operaciones realizadas para que se pueda ver
f�cilmente si se producen errores en el sistema.


\subsection{Lista de requisitos funcionales}

Despu�s de analizar con detalle los requisitos y las consideraciones tenidas
en cuenta para el desarrollo del sistema, en este apartado se listan todos los
requisitos funcionales del sistema.

\subsubsection{Servidor front-end}

\begin{milista}
\item \textbf{RF1}: Identificar y registrar a los clientes.
\item \textbf{RF2}: Obtener los datos de un beneficiario.
\item \textbf{RF3}: Almacenar un nuevo beneficiario.
\item \textbf{RF4}: Modificar los datos de un beneficiario.
\item \textbf{RF5}: Eliminar un beneficiario.
\item \textbf{RF6}: Obtener los datos de un usuario (o m�dico).
\item \textbf{RF7}: Almacenar un nuevo usuario (o m�dico).
\item \textbf{RF8}: Modificar los datos de un usuario (o m�dico).
\item \textbf{RF9}: Eliminar un usuario (o m�dico).
\item \textbf{RF10}: Almacenar sustituciones a m�dicos.
\item \textbf{RF11}: Consultar las citas de un beneficiario.
\item \textbf{RF12}: Consultar las citas de un m�dico.
\item \textbf{RF13}: Almacenar una nueva cita.
\item \textbf{RF14}: Eliminar una cita.
\item \textbf{RF15}: Almacenar un nuevo volante.
\item \textbf{RF16}: Registrar todas las operaciones realizadas por los
usuarios en la BD.
\item \textbf{RF17}: Actualizar el estado de los clientes cuando se produce una
operaci�n en el servidor.
\item \textbf{RF18}: Enviar peticiones al servidor de respaldo para gestionar
la base de datos secundaria.
\end{milista}

\subsubsection{Servidor de respaldo}

\begin{milista}
\item \textbf{RF1}: Recibir peticiones para gestionar la base de datos secundaria.
\end{milista}

\subsubsection{Unidad de citaci�n}

\begin{milista}
\item \textbf{RF1}: Identificar a los usuarios del sistema.
\item \textbf{RF2}: Consultar los datos de un beneficiario.
\item \textbf{RF3}: Registrar un nuevo beneficiario.
\item \textbf{RF4}: Modificar los datos de un beneficiario.
\item \textbf{RF5}: Eliminar un beneficiario.
\item \textbf{RF6}: Consultar los datos de un usuario (o m�dico).
\item \textbf{RF7}: Registrar un nuevo usuario (o m�dico).
\item \textbf{RF8}: Modificar los datos de un usuario (o m�dico).
\item \textbf{RF9}: Eliminar un usuario (o m�dico).
\item \textbf{RF10}: Asignar sustitutos a los m�dicos.
\item \textbf{RF11}: Consultar las citas de un beneficiario.
\item \textbf{RF12}: Consultar las citas de un m�dico.
\item \textbf{RF13}: Registrar una nueva cita.
\item \textbf{RF14}: Anular una cita.
\item \textbf{RF15}: Emitir un volante.
\item \textbf{RF16}: Recibir notificaciones para actualizar el estado.
\end{milista}


\clearpage

% Secci�n de ap�ndices
% Se a�ade el titulo de "Ap�ndices" al �ndice
\addcontentsline{toc}{section}{Ap�ndices}

%\input{Apendices.tex}

% A�adimos la bibliografia al �ndice
\clearpage\phantomsection
\addcontentsline{toc}{section}{\bibname}

% Bibliograf�a. En este caso se usa BibTeX
%\bibliographystyle{plain}
%\pagestyle{plain} 
%\bibliography{Bibliografia}

\end{document}
